\documentclass[10pt]{article}
\usepackage[preprint]{tmlr}
\makeatletter
\@acceptedtrue
\@preprinttrue
\makeatother
\input{math_commands.tex}

\usepackage{graphicx}
\usepackage{booktabs}
\usepackage{tabularx}
\usepackage{array}
\usepackage{amsmath}
\usepackage{amssymb}
\usepackage{hyperref}
\usepackage{url}
\usepackage{microtype}
\usepackage{needspace}
\hypersetup{hidelinks}

\title{A Reproducible Benchmark of Imbalance Handling\\for Multiclass Tabular Classification}

% De-anonymized local draft. Confirm affiliation and venue-specific metadata
% before any external submission.
\author{Shaurya Saria\\\texttt{sariashaurya09@gmail.com}}

\newcommand{\method}[1]{\texttt{#1}}
\newcommand{\metric}[1]{\textit{#1}}
\newcommand{\analysisfigure}[3]{%
  \begin{figure}[t]
    \centering
    \IfFileExists{#1}{%
      \includegraphics[width=\linewidth]{#1}%
    }{%
      \fbox{\parbox{0.88\linewidth}{\centering
        The figure is generated by the locked analysis pipeline.\\
        Expected local path: \texttt{#1}}}%
    }
    \caption{#2}
    \label{fig:#3}
  \end{figure}%
}

\begin{document}

\maketitle

\begin{abstract}
Class imbalance in multiclass tabular classification is commonly studied with incomplete comparisons and fold-level inference. We present a preregistered, reproducible benchmark of 11 public datasets (9 numeric, 2 mixed), three classifiers, and locked weighting and sampling conditions. Preprocessing is fit within training folds; evaluation uses stratified five-fold cross-validation under three seeds, dataset-level complete-cell aggregation, macro-F1 as the primary outcome, and paired uncertainty-aware comparisons. The run produced 4,191 valid and 1,254 failed records from 5,445 expected records. Four of 32 Friedman tests were significant, all for numeric XGBoost; Nemenyi tests isolated SMOTEENN versus SMOTETomek. Among 184 Holm-adjusted raw-baseline comparisons, one was significant: SMOTEENN reduced MCC for numeric logistic regression. These results do not support a universally best sampler. They instead show that benchmark conclusions depend on applicability, complete-cell coverage, paired effects, and computational cost. Claims are limited to the locked 11-dataset scope and observed coverage.
\end{abstract}

\section{Introduction}

Class imbalance is a recurring source of difficulty in multiclass classification. A model can obtain a strong aggregate score while failing systematically on a rare class, and a resampling method that helps one classifier or feature representation can be neutral or harmful for another. This is especially important for tabular data, where the interaction between preprocessing, categorical representation, model inductive bias, and the sampling operator is part of the empirical question rather than a disposable implementation detail.

The literature contains mature families of imbalance interventions, including class weighting, random over- and under-sampling, synthetic minority generation, and combinations of over-sampling with cleaning or boundary-based methods \citep{he2009learning,branco2016survey}. Their comparison is nevertheless difficult to interpret when datasets are selected after pilot scores, unsupported combinations are silently dropped, folds are treated as independent replications, or one metric conceals class-specific failures. Benchmark studies in adjacent tabular settings show the value of fixed data collections, explicit protocols, and broad model comparisons \citep{gorishniy2022why}; the same principles matter when preprocessing and sampling choices interact.

We ask how locked conditions compare with raw baselines; whether results differ by feature family and classifier; and what coverage, failure, runtime, and post-sampling row-count costs accompany the metric changes. Scope, conditions, metric definitions, and complete-case rules were fixed before analysis.

Our contributions are threefold:

\begin{itemize}
  \item A locked comparison of eleven public datasets, three classifiers, and feature-type-appropriate conditions.
  \item A fold-safe protocol that retains unsupported and failed runs as explicit audit rows.
  \item Dataset-level paired effects, uncertainty, statistical tests, runtime, row counts, and reproducible source.
\end{itemize}

The resulting evidence supports a narrow conclusion: imbalance handling is conditional on the dataset, classifier, feature family, and metric, and the cost of a method is part of its practical evaluation. The paper does not claim a universally superior sampler.

\section{Related Work}

\paragraph{Imbalance interventions.}
SMOTE generates synthetic minority examples by interpolating between minority neighbours \citep{chawla2002smote}; ADASYN adapts the amount of synthesis to local learning difficulty \citep{he2008adasyn}; and Borderline-SMOTE concentrates synthesis near difficult class boundaries \citep{han2005borderline}. Random over- and under-sampling provide simpler distributional changes. SMOTEENN and SMOTETomek combine over-sampling with cleaning or Tomek-link removal, respectively, and are representative of combination approaches studied in the imbalance literature \citep{batista2004study}. Our benchmark treats these methods as locked comparators, not as tunable algorithmic contributions.

\paragraph{Models and metrics.}
The model set intentionally spans a linear classifier, a bagged tree ensemble, and gradient-boosted trees. Balanced Random Forest is included as an RF-specific comparator following the idea of drawing balanced bootstrap samples for trees \citep{chen2004using}. XGBoost is a widely used gradient-boosting implementation for tabular prediction \citep{chen2016xgboost}. We use macro-F1 as the primary metric and report multiclass G-mean, Matthews correlation coefficient, balanced accuracy, and per-class recall as secondary outcomes. This combination reflects the need to evaluate both class-wise performance and agreement beyond overall accuracy \citep{sokolova2009systematic,gorodkin2004comparing}.

\paragraph{Benchmarking and reproducibility.}
Benchmark conclusions depend on the data collection, comparison matrix, aggregation level, and statistical block definition. Recent tabular benchmark work demonstrates how a carefully controlled collection can reveal broad model patterns without making every dataset or model appear interchangeable \citep{gorishniy2022why}. The present benchmark narrows that principle to imbalance handling: all condition applicability decisions are protocol-defined, dataset-level means are the inferential observations, and the raw failure table remains part of the evidence. The implementation uses the open-source imbalanced-learn ecosystem for the sampling operators \citep{lemaitre2017imbalanced}.

\section{Benchmark Design}
\label{sec:design}

\subsection{Datasets and scope}

The benchmark was scoped to independent, public, multiclass tabular datasets with at least three classes, documented source terms, and sufficient support for five-fold stratification under the locked feasibility rule. The candidate registry contained 25 OpenML/UCI records. Eligibility and feature-type decisions were made from provenance, target, leakage, duplicate, missingness, and support audits; preliminary scores were not used to retain or remove datasets. Eleven datasets were retained: nine numeric datasets and two mixed datasets (CMC and Dermatology). Multiple Features (Factors) was deferred because its corrected replay took 1,043.5 seconds against the 900-second Stage 0 feasibility budget. This is a computational scope decision, not a score-based exclusion.

Table~\ref{tab:datasets} reports the retained registry fields used in the analysis. The imbalance ratio is the majority-class count divided by the minority-class count. The retained collection spans 150--5,620 rows, 3--10 classes, 4--856 raw predictors, and imbalance ratios from 1.0 to 92.6. Minimum class support ranges from 5 to 554, making the complete-cell and failure rules important for interpreting the result coverage.

\begin{table}[t]
\centering
\caption{Retained datasets and registry-level moderators.}
\label{tab:datasets}
\small
\begin{tabular}{@{}lrrrrlr@{}}
\toprule
Dataset & $n$ & $C$ & $d$ & $n_{\min}$ & Type & IR \\
\midrule
Balance Scale & 625 & 3 & 4 & 49 & Numeric & 5.88 \\
CNAE-9 & 1,080 & 9 & 856 & 120 & Numeric & 1.00 \\
Glass & 214 & 6 & 9 & 9 & Numeric & 8.44 \\
Iris & 150 & 3 & 4 & 50 & Numeric & 1.00 \\
Optdigits & 5,620 & 10 & 64 & 554 & Numeric & 1.03 \\
Seeds & 210 & 3 & 7 & 70 & Numeric & 1.00 \\
Wine & 178 & 3 & 13 & 48 & Numeric & 1.48 \\
Wine Quality Red & 1,599 & 6 & 11 & 10 & Numeric & 68.10 \\
Yeast & 1,484 & 10 & 8 & 5 & Numeric & 92.60 \\
CMC & 1,473 & 3 & 9 & 333 & Mixed & 1.89 \\
Dermatology & 366 & 6 & 34 & 20 & Mixed & 5.60 \\
\bottomrule
\end{tabular}
\end{table}

The raw files, source identifiers, retrieval metadata, licenses or terms, target semantics, class counts, and SHA-256 hashes are recorded in the versioned dataset registry and acquisition manifest. Raw and processed data are not committed to the source repository; the manifests are the provenance boundary for reproducing the local evidence package.

\subsection{Evaluation protocol}

Each dataset uses stratified five-fold cross-validation with shuffling and seeds $0$, $1$, and $2$. A required record is identified by dataset, classifier, condition, seed, and fold, yielding 15 seed-fold records per cell. Rows with missing targets are removed before splitting. Predictor imputation, encoding, scaling, resampling, and training-fold weights are fitted within each training fold. The held-out fold is only transformed and predicted; it is never resampled or used to fit a preprocessing operation.

Numeric pipelines impute numeric values with the training median and scale after imputation. Mixed pipelines impute numeric values with the training median and categorical values with the training mode, ordinal-encode categorical predictors before SMOTENC, and one-hot encode after resampling with unknown categories ignored. Semantic categorical overrides are recorded separately from the registry so that numeric source codes are not treated as continuous by accident.

\subsection{Models and conditions}

The classifiers are logistic regression, random forest, and XGBoost. Logistic regression uses a maximum of 1,000 iterations. Random forest uses 200 trees and one worker. XGBoost uses 200 estimators, the histogram tree method, multiclass log loss for evaluation, and one worker. The locked conditions are summarized in Table~\ref{tab:conditions}; sampler parameters are deliberately simple and fixed rather than tuned after seeing the full-run scores.

\begin{table}[t]
\centering
\caption{Locked comparison matrix and applicability rules.}
\label{tab:conditions}
\small
\begin{tabularx}{\linewidth}{@{}l>{\raggedright\arraybackslash}X>{\raggedright\arraybackslash}p{0.24\linewidth}@{}}
\toprule
Family & Conditions & Applicability \\
\midrule
Baseline/weighting & Raw; class-weighted with the balanced class-weight rule & All feature families; model-specific class weights \\
Simple resampling & Random over-sampling; random under-sampling & All feature families \\
Numeric synthesis & SMOTE; ADASYN; Borderline-SMOTE (borderline-1) & Numeric datasets; three neighbours \\
Mixed synthesis & SMOTENC; SMOTEENN; SMOTETomek & Mixed datasets; ordinal representation before resampling \\
Balanced ensemble & Balanced Random Forest & Random Forest only; 200 trees; replacement \\
\bottomrule
\end{tabularx}
\end{table}

The numeric core contains nine conditions for each classifier: raw, class-weighted, random over-sampling, random under-sampling, SMOTE, ADASYN, Borderline-SMOTE, SMOTEENN, and SMOTETomek. The mixed core contains seven conditions, replacing the ordinary synthetic methods with SMOTENC, SMOTEENN, and SMOTETomek. Balanced Random Forest is analysed only within the RF classifier and is kept separate from the general classifier comparisons.

\subsection{Metrics and aggregation}

Macro-F1 is the primary outcome, computed over all declared target classes with zero division set to zero. Secondary outcomes are multiclass G-mean, MCC, balanced accuracy, and per-class recall. If $R_c$ denotes recall for class $c$, then

\begin{equation}
G = \left(\prod_{c=1}^{C} R_c\right)^{1/C}.
\label{eq:gmean}
\end{equation}

All encoded classes are included in Equation~\ref{eq:gmean}; a zero class recall makes G-mean exactly zero. MCC is the multiclass Matthews coefficient, balanced accuracy is the unweighted mean recall, and per-class recall uses the same declared labels and zero behavior.

Metric values are first averaged over the 15 seed-fold records within a complete dataset--classifier--condition cell. A cell is complete only if all 15 records are valid. The statistical block is the dataset, not the fold or seed. For each family, classifier, and metric, only datasets with a complete cell for every condition in that family are included. This produces paired dataset-level comparisons and prevents a method with more failures from appearing artificially strong because only its easiest cells remain.

\subsection{Statistics and efficiency}

For every non-baseline condition, the paired difference is the condition mean minus the raw-baseline mean for the same dataset, classifier, and metric. We report mean and median paired differences, win/tie/loss counts, paired rank-biserial correlation, and a 95\% percentile bootstrap interval with 10,000 resamples and NumPy generator seed 0. Ties are defined using the locked $10^{-12}$ tolerance.

The confirmatory block uses two-sided tests with $\alpha=0.05$. Friedman tests \citep{friedman1937use} compare all conditions in each complete family when at least three dataset blocks and three conditions are available. Significant Friedman tests are followed by the pre-specified all-pairs Nemenyi procedure \citep{nemenyi1963distribution}. Limited raw-versus-alternative comparisons use two-sided paired Wilcoxon signed-rank tests \citep{wilcoxon1945individual} with zero differences discarded; Holm correction is applied within each family, classifier, and metric. Macro-F1 is the sole primary confirmatory metric; the secondary metrics are interpreted as supporting evidence. A comparison with fewer than three complete blocks is retained as a skipped or descriptive result with its reason.

Runtime, peak resident memory, and training rows before and after sampling are descriptive efficiency outcomes. Imbalance ratio, class count, raw and encoded dimensionality, and feature type are exploratory moderators. With eleven datasets, moderator patterns are descriptive and cannot establish a selection rule.

\subsection{Reproducibility boundary}

The source package freezes the configuration, registry, acquisition manifest, environment record, scope lock, raw-run manifest, analysis manifest, and Gate E review. The generated raw and analysis results remain ignored by Git but are hashed in the local release payload. The executable scripts and tests preserve the fold boundaries, validity accounting, aggregation, statistical rules, and figure generation. The paper is a draft for review; it does not imply external submission or venue acceptance.

\section{Results}
\label{sec:results}

\subsection{Coverage and failures}

The full run enumerated 5,445 expected records, of which 4,191 were valid and 1,254 were explicit failures. The analysis package contains 278 complete dataset--classifier--condition cells. Complete-cell counts were 75, 84, and 75 for logistic regression, random forest, and XGBoost in the numeric family, respectively; the corresponding mixed-family counts were 14, 16, and 14. The additional RF counts reflect the Balanced Random Forest comparator. These counts are not interchangeable with the number of folds: every complete cell contains 15 seed-fold records.

Failures are interpretable protocol and applicability evidence. The grouped failure messages contain 405 records for applying SMOTENC outside mixed data, 330 records for applying Balanced Random Forest outside RF, 270 records for ordinary numeric synthetic samplers declared invalid for mixed data in Stage 0, 207 records in which the requested ratio generated no samples, and 42 ADASYN records reporting that no majority neighbours were available for the dataset. Unsupported combinations and sampler-specific failures therefore remain visible in the denominator rather than being treated as poor scores or silently discarded.

\analysisfigure{../results/analysis/figures/class_distribution_overview.png}{Class distributions for the 11 retained datasets; descriptive evidence for the moderator range.}{class-distribution}

\subsection{Descriptive performance}

Figure~\ref{fig:macro-f1-heatmap} shows the complete-cell macro-F1 means by classifier, condition, and dataset. Because each condition can have a different complete-cell set, the heatmap is descriptive and should not be read as a single pooled estimate. In the numeric family, the complete-cell mean for ADASYN is high for all three classifiers (0.9465 for logistic regression, 0.9332 for random forest, and 0.9242 for XGBoost), but only three numeric dataset blocks are complete for that condition. Its paired baseline comparison consequently has only three pairs, all ties. The result is evidence of limited applicability in this run, not evidence that ADASYN is generally superior.

For the nine-pair numeric comparisons, the largest observed mean paired macro-F1 difference was +0.0306 for logistic regression with random over-sampling, +0.0074 for random forest with SMOTETomek, and +0.0076 for XGBoost with class weighting. Their bootstrap intervals include zero for logistic regression and are narrow but not Holm-significant for the two tree-based comparisons. In the mixed family, the largest observed mean paired differences were +0.0127 for logistic regression with class weighting, +0.0149 for random forest with SMOTEENN, and +0.0113 for XGBoost with SMOTEENN; each uses only two mixed dataset blocks and is therefore descriptive.

\begin{table}[t]
\centering
\caption{Largest mean paired macro-F1 differences from the raw baseline by family and classifier. Intervals are bootstrap confidence intervals; rows are descriptive, not score-selected.}
\label{tab:macro-deltas}
\scriptsize
\begin{tabular}{@{}lllrrrrr@{}}
\toprule
Family & Classifier & Condition & $n$ & Mean $\Delta$ & 95\% CI & W/T/L & Holm $p$ \\
\midrule
Numeric & LR & Random over & 9 & +0.0306 & [-0.0134, 0.0930] & 3/3/3 & 1.000 \\
Numeric & RF & SMOTETomek & 9 & +0.0074 & [0.0016, 0.0148] & 7/0/2 & 0.313 \\
Numeric & XGBoost & Class weighted & 9 & +0.0076 & [0.0015, 0.0150] & 6/3/0 & 0.219 \\
Mixed & LR & Class weighted & 2 & +0.0127 & [0.0010, 0.0244] & 2/0/0 & 1.000 \\
Mixed & RF & SMOTEENN & 2 & +0.0149 & [-0.0059, 0.0356] & 1/0/1 & 1.000 \\
Mixed & XGBoost & SMOTEENN & 2 & +0.0113 & [-0.0076, 0.0303] & 1/0/1 & 1.000 \\
\bottomrule
\end{tabular}
\end{table}

\analysisfigure{../results/analysis/figures/macro_f1_ranking_heatmap.png}{Complete-cell macro-F1 means by classifier, condition, and dataset. Rankings are descriptive; inference uses paired complete cases.}{macro-f1-heatmap}

\subsection{Confirmatory analysis}

There were 32 planned Friedman rows. Four were significant, all in the numeric-core XGBoost block: macro-F1, G-mean, MCC, and balanced accuracy each had Friedman statistic 24.0 and $p=0.002292$ across three complete dataset blocks and nine conditions. The associated Nemenyi results identified SMOTEENN versus SMOTETomek as the only $p<0.05$ pair for each of those four metrics, with $p=0.010464$. The remaining 28 Friedman rows were non-significant or skipped. In particular, mixed-family tests had only two complete dataset blocks and were reported as skipped rather than over-interpreted.

Across 184 raw-versus-alternative Wilcoxon rows, exactly one remained significant after Holm correction. For numeric logistic regression and MCC, SMOTEENN had mean paired delta $-0.0667$, median delta $-0.0265$, zero wins, zero ties, nine losses, rank-biserial correlation $-1.0$, and a 95\% bootstrap interval of $[-0.1105,-0.0283]$; the unadjusted $p$-value was 0.003906 and the Holm-adjusted value was 0.03125. This is a statistically supported negative comparison, not a positive sampler recommendation. The remaining corrected baseline comparisons did not meet the declared threshold.

\analysisfigure{../results/analysis/figures/macro_f1_baseline_deltas.png}{Paired macro-F1 differences from the raw baseline. Positive values indicate improvement; intervals and pair counts are retained in the analysis output.}{macro-f1-deltas}

The class-wise results are reported in the analysis package rather than collapsed into one cross-dataset class index, because class labels and class semantics differ across datasets. The per-class recall summaries provide descriptive evidence for which classes drive a macro metric while preserving the dataset-specific label space.

\Needspace{9\baselineskip}
\subsection{Efficiency costs}

Sampling changed the number of training rows in the expected direction, but the magnitude depended strongly on the dataset. Across raw numeric cells, the average training set size before sampling was 992 rows. SMOTE increased this to an average of 1,548 rows, random over-sampling had the same average post-sampling size, SMOTETomek averaged 1,518 rows, and SMOTEENN averaged 1,247 rows; random under-sampling reduced it to 661 rows. These are descriptive averages over dataset-level efficiency rows, not a claim that every dataset follows the same pattern.

The average fit-and-predict time for numeric random forest was 0.474 seconds for the raw baseline and 0.625 seconds for SMOTETomek; Balanced Random Forest averaged 0.721 seconds. For numeric XGBoost, the raw baseline averaged 0.538 seconds, while SMOTETomek averaged 0.737 seconds. The numeric ADASYN mean for XGBoost was 0.860 seconds, again based on only three complete datasets. In the two mixed datasets, RF averaged 0.365 seconds for the raw baseline, 0.501 seconds for SMOTENC, and 0.538 seconds for Balanced Random Forest. These measurements support treating computational cost and applicability as first-class benchmark outcomes.

\section{Discussion}
\label{sec:discussion}

The main result is conditionality rather than a leaderboard. A few numeric XGBoost comparisons separated in the all-condition Friedman/Nemenyi block, while the only Holm-significant raw-baseline comparison was a negative MCC result for SMOTEENN with logistic regression. Several descriptive macro-F1 deltas were positive, but their intervals or corrected $p$-values do not support a general recommendation. The high ADASYN complete-cell means illustrate the same point from the coverage side: the method is complete on only three numeric dataset blocks, and its paired comparison has no non-zero differences. Reporting the coverage and failure denominator changes the interpretation of the raw mean.

\Needspace{9\baselineskip}
The benchmark also shows why feature type should be part of the protocol. Ordinary SMOTE-family conditions are not applied to mixed one-hot features; mixed data uses SMOTENC before one-hot encoding. As a result, the mixed family has a different comparison matrix and only two retained dataset blocks. Pooling numeric and mixed conditions would either violate applicability or conceal that the mixed inference block is too small for the planned tests.

\Needspace{8\baselineskip}
The efficiency evidence reinforces this interpretation. Over-sampling usually increased the number of training rows, and more elaborate conditions generally increased fit-and-predict time relative to the raw baseline. The effect is not a universal penalty: random under-sampling reduced row counts and often reduced runtime, but its performance deltas varied by classifier and dataset. A practical choice therefore involves a task-specific trade-off among class-wise performance, failure risk, resource cost, and reproducibility.

\subsection{Limitations}

First, eleven datasets are enough for a controlled benchmark but not enough to establish a universal method ranking, especially for the two-dataset mixed family. Second, the candidate pool and retained scope are fixed by the project protocol and are not a claim to represent all tabular domains. Third, sampler and classifier hyperparameters are locked baseline settings rather than a full tuning study; tuning could change both performance and cost. Fourth, the complete-cell rule is conservative but reduces the inferential sample when a method fails or is inapplicable. Fifth, exploratory moderators are descriptive and should not be used as a post-hoc selection rule. Finally, the current artifact is a technical paper draft and local evidence package, not an externally reviewed or venue-submitted manuscript.

\subsection{Reproducibility and future work}

The repository contains the protocol, registry, acquisition provenance, executable benchmark and analysis scripts, focused tests, generated-figure definitions, and Gate F/G snapshot validators. A clean-checkout reproduction revalidates the source and frozen manifests without silently widening the experiment. Future work could enlarge the retained dataset collection, implement categorical-only handling, add carefully pre-registered tuning budgets, and test whether the observed conditionality persists under other model families. Such extensions would require a new scope or protocol amendment rather than being folded into this result set.

\section{Conclusion}

We presented a reproducible benchmark of imbalance handling for multiclass tabular classification. Under a locked eleven-dataset protocol, explicit applicability and failure accounting, dataset-level complete-cell aggregation, and uncertainty-aware paired analyses, the evidence did not identify a universally best condition. The most defensible conclusion is methodological: imbalance benchmarks should report what ran, what failed, how much data and time each condition consumed, and how much evidence remains after complete-case filtering. Those reporting choices make modest and negative results useful rather than misleading.

\section*{Acknowledgments}

The author thanks the maintainers and contributors of the open-source
scientific Python ecosystem, particularly scikit-learn and imbalanced-learn,
whose software supported this benchmark.

\bibliography{references}
\bibliographystyle{tmlr}

\end{document}
